\section{Explicit reconstruction from complete returns}
\label{ret:seccion}

\subsection{State domain and bilateral extension}

Consider a chart of the surviving domain generated by
APP--TRIT--TPK. A state in this chart preserves the operator sheets,
the tritic orientation, the remainder, the quotient, the carry, the route,
the boundary, the memory record, and the schedule origin.
The update acts on this complete state.

Let \(X_0\) be a nonempty compact Hausdorff space and
\(F:X_0\to X_0\) a continuous surjective update.
These are the hypotheses of the result in the chart under consideration.
The bilateral extension of the update is
\begin{equation}
 \mathscr X_F=
 \bigl\{(x_k)_{k\in\ZZ}\in X_0^{\ZZ}:
                  F(x_k)=x_{k+1}\text{ for every }k\in\ZZ\bigr\},
 \label{ret:extension}
\end{equation}
with the topology induced by the product.

\begin{lemma}[Existence and bilateral transport]
\label{ret:existencia}
The space \(\mathscr X_F\) is compact, Hausdorff, and nonempty.
Every coordinate projection \(\mathscr X_F\to X_0\) is surjective.
The shift
\[
 \sigma((x_k)_{k\in\ZZ})=(x_{k+1})_{k\in\ZZ}
\]
is a homeomorphism of \(\mathscr X_F\).
\end{lemma}
\begin{proof}
Each condition \(F(x_k)=x_{k+1}\) defines a closed subset of the compact
Hausdorff product \(X_0^{\ZZ}\).
For a fixed index \(k_0\) and state \(z\in X_0\), any finite
collection of these conditions, together with \(x_{k_0}=z\), admits a solution:
complete forwards using \(F\) and backwards using
preimages, which exist by surjectivity. The conditions therefore have
the finite intersection property. Compactness produces
a bilateral history with \(x_{k_0}=z\).

The shifts \(\sigma\) and
\(\sigma^{-1}((x_k))=(x_{k-1})\) preserve the relations in
\eqref{ret:extension}; both are continuous by their coordinate action.
\end{proof}

\subsection{Tritic sampling and reconstruction of every step}

With the schedule origin fixed, the typed return times are
\[
 \mathcal T_+
 =\{k\in\ZZ:1+(k\bmod 9)\in\{1,4,7\}\}=3\ZZ.
\]
The complete state \(x_k\) is preserved at each such time.
Define the space of return records by
\begin{equation}
 \mathscr Y_+
 =\bigl\{(y_j)_{j\in\ZZ}\in X_0^{\ZZ}:
                 F^3(y_j)=y_{j+1}\text{ for every }j\in\ZZ\bigr\}.
 \label{ret:registros}
\end{equation}

\begin{theorem}[Recovery from complete returns]
\label{ret:homeomorfismo}
The map
\[
 r=\operatorname{Ret}_+:\mathscr X_F\longrightarrow\mathscr Y_+,
 \qquad r(x)_j=x_{3j},
\]
is a homeomorphism. Its inverse is given explicitly by
\begin{equation}
 \boxed{
  \bigl(r^{-1}(y)\bigr)_{3j+s}=F^s(y_j),
  \qquad j\in\ZZ,\quad s\in\{0,1,2\}.}
 \label{ret:inversa}
\end{equation}
\end{theorem}
\begin{proof}
For \(x\in\mathscr X_F\),
\(F^3(x_{3j})=x_{3j+3}\), so \(r(x)\in\mathscr Y_+\).
Conversely, every integer \(k\) has a unique expression
\(k=3j+s\), with \(j\in\ZZ\) and \(0\le s<3\), including when
\(k<0\). Formula \eqref{ret:inversa} therefore defines every
coordinate of a sequence \(x\).

If \(s=0\) or \(s=1\), then
\(F(x_{3j+s})=F^{s+1}(y_j)=x_{3j+s+1}\).
At the boundary between two blocks,
\[
 F(x_{3j+2})=F^3(y_j)=y_{j+1}=x_{3j+3}.
\]
Thus \(x\in\mathscr X_F\).
Reading its coordinates indexed by \(3j\) returns \(y_j\).
In the reverse direction, the relations in \eqref{ret:extension} give
\(x_{3j+s}=F^s(x_{3j})\), so reconstruction from \(r(x)\)
returns \(x\) in full.

The map \(r\) is continuous because it selects coordinates.
Each coordinate of its inverse is the composition of a coordinate
projection with one of the continuous maps \(I,F,F^2\).
The inverse is therefore continuous.
\end{proof}

An additional discrete coordinate, such as the index \(m\in\ZZ\)
of a realization chart, is preserved through
\[
 \widetilde r:\ZZ\times\mathscr X_F
 \longrightarrow\ZZ\times\mathscr Y_+,
 \qquad \widetilde r(m,x)=(m,r(x)).
\]
Its inverse is \((m,y)\mapsto(m,r^{-1}(y))\).
The same result holds for a fixed schedule origin \(a\):
replace \(x_{3j}\) and \(x_{3j+s}\) with
\(x_{a+3j}\) and \(x_{a+3j+s}\).

\subsection{Conjugation of transport and preservation of readouts}

Let \(\sigma_+(y)_j=y_{j+1}\) be the shift of records.
The preceding reconstruction gives
\begin{equation}
 r\sigma^3=\sigma_+r,\qquad
 (r\sigma r^{-1}(y))_j=F(y_j).
 \label{ret:conjugacion}
\end{equation}
The map in the second equality is reversible on
\(\mathscr Y_+\); its inverse is
\[
 y\longmapsto\bigl(F^2(y_{j-1})\bigr)_{j\in\ZZ}.
\]
Indeed, both compositions reduce to
\(F^3(y_{j-1})=y_j\).
Reversibility of the bilateral history is thus preserved even
when the update \(F\) on an isolated state is not injective.
If the nonadic revolution is realized as \(\Gamma_9=\sigma^9\), then
\begin{equation}
 r\Gamma_9=\sigma_+^3r.
 \label{ret:vuelta-nonadica}
\end{equation}

\begin{corollary}[Joint readouts and recovered transport]
\label{ret:lecturas}
Let \(a:\mathscr X_F\to\mathcal R\) and
\(q:\mathscr X_F\to\mathcal E\) be two continuous readers of the same
state, and let \(d:\mathscr X_F\to\mathscr X_F\) be a homeomorphism of the
calibrated domain. Then
\[
 a_+=a\circ r^{-1},\qquad
 q_+=q\circ r^{-1},\qquad
 d_+=r\circ d\circ r^{-1}
\]
are continuous, \(d_+\) is a homeomorphism, and
\begin{equation}
 (a,q)=(a_+,q_+)\circ r,\qquad rd=d_+r.
 \label{ret:factorizacion}
\end{equation}
For every integer \(N>0\),
\[
 d^Nx=x\quad\Longleftrightarrow\quad d_+^Nr(x)=r(x).
\]
\end{corollary}
\begin{proof}
Compose the maps with the homeomorphism of
\cref{ret:homeomorfismo}. The inverse of \(d_+\) is
\(rd^{-1}r^{-1}\). The identity
\(d_+^Nr=rd^N\), together with injectivity of \(r\), proves the
last equivalence.
\end{proof}

Application to the double readout takes
\(\mathcal R=R_{\rm HMT}\) and the incidence codomain constructed in
\cref{exc:doble-lectura,exc:limite-inverso}.
The result establishes a recoverable presentation of the same
histories, their readers, and their transport. Its object is this
dynamical equivalence of presentations.
